\section{Conclusiones}

El proyecto cumplió su objetivo general: evolucionar la arquitectura de Paralegales hacia un diseño híbrido Firebase y AWS sin disrumpir la operación de una plataforma en producción con base de usuarios activa. Esta última condición —madurar el sistema sin detener el servicio— fue la restricción más exigente del trabajo y orientó la decisión de fondo: optar por una evolución arquitectónica incremental y controlada en lugar de una migración total disruptiva.

Los cuatro frentes complementarios que articularon esa evolución contribuyeron, cada uno, a la maduración técnica de la plataforma. La reubicación de la lógica de negocio crítica hacia entornos del lado del servidor mediante el patrón \textit{Backend-for-Frontend} redujo la superficie de exposición y habilitó la instrumentación de los flujos sensibles. La refactorización modular —Nuxt Layers bajo \textit{Domain-Driven Design} en el \textit{frontend} y la separación \textit{Handler--Service--Repository} en el \textit{backend}— disminuyó el acoplamiento y habilitó la evolución independiente de cada dominio y de cada motor tecnológico subyacente. La observabilidad estructurada, con identificadores de correlación e infraestructura definida como código, otorgó visibilidad del comportamiento del sistema en producción a través de componentes heterogéneos. Y las prácticas automatizadas de aseguramiento de calidad permitieron sostener la velocidad de cambio sin degradar la confiabilidad ni la seguridad del sistema.

Más allá de los artefactos construidos, el trabajo deja una reflexión de ingeniería: la velocidad inicial que habilitó Firebase, valiosa para validar el producto, debía equilibrarse con la sostenibilidad a largo plazo, y ese equilibrio no exigió rehacer el sistema, sino evolucionarlo de forma disciplinada. La estrategia de evolución arquitectónica controlada se confirma así como un camino viable para madurar un sistema real sin interrumpir su operación.

En conjunto, el diseño resultante deja a la plataforma con menor exposición de su lógica crítica, mayor trazabilidad operativa y mejor mantenibilidad, preparada para una evolución incremental continua a medida que el sistema y su base de usuarios crecen.
